\subsection{Introduction to Relations} 
\begin{definition}[Relation]
A \term{relation} is a set of ordered pairs. If $\relR$ is a relation, we write $x\relR y$ to mean that the pair $(x,y)$ is an element of the relation $\relR$.
\end{definition}
For our purposes, when we work with a relation, each pair will be a pair of \emph{numbers}. Unless we specify otherwise, it will be a pair of \makeblank[1in] numbers.
\begin{Example}
Define a relation by
\[
R=\left\{(1,3),(3,5),(5,-1)\right\}
\]
\begin{itemize}
\item The statement $1\relR 3$ is \makeblank[1in] because $(1,3)$ \makeblank[1in] an element of $\relR$
\item The statement $4\relR 2$ is \makeblank[1in] because $(2,4)$ \makeblank[1in] an element of $\relR$
\item The statement $5\relR 3$ is \makeblank[1in] because $(5,3)$ \makeblank[1in] an element of $\relR$
\end{itemize}
\end{Example}
We can also define a relation in \makeblank notation.
\begin{Example}
Define a relation by
\[
S=\left\{(x,y)\mid x+y=10\right\}
\]
To determine whether $x\relS y$ is true, we must determine whether \blankpair\ is an element of $S$, that is, whether \makeblank[1in] is true.
\begin{itemize}
\item $2\relS 8$ is \makeblank[1in], since \makeblank[1.5in].
\item $5\relS 1$ is \makeblank[1in], since \makeblank[1.5in].
\item $19\relS -9$ is \makeblank[1in], since \makeblank[1.5in].
\end{itemize}
\end{Example}
\begin{Example}
Define a relation by
\[
T=\left\{(x,y)\mid x<y\right\}
\]
To determine whether $x\relT y$ is true, we must determine whether \blankpair\ is an element of $T$, that is, whether \makeblank[1in] is true.
\begin{itemize}
\item $-1\relT 6$ is \makeblank[1in], since \makeblank[1in] is \makeblank[1in].
\item $7\relT 4$ is \makeblank[1in], since \makeblank[1in] is \makeblank[1in].
\item $3\relT 3$ is \makeblank[1in], since \makeblank[1in] is \makeblank[1in].
\end{itemize}
\end{Example}